\documentclass[11pt]{article}
\usepackage[margin=1in]{geometry}
\usepackage{amsmath,amssymb,amsthm}
\usepackage{array}
\usepackage{parskip}
\usepackage{xeCJK}

\newtheorem{theorem}{Theorem}
\newtheorem{lemma}[theorem]{Lemma}
\newtheorem{corollary}[theorem]{Corollary}
\newtheorem{proposition}[theorem]{Proposition}
\theoremstyle{definition}
\newtheorem{definition}[theorem]{Definition}
\theoremstyle{remark}
\newtheorem{remark}[theorem]{Remark}

\newcommand{\N}{\mathbb{N}}
\DeclareMathOperator{\cont}{content}
\DeclareMathOperator{\SYT}{SYT}
\DeclareMathOperator{\SSYT}{SSYT}

\title{Disproof of conjecture \texttt{00000000418}}
\author{earthking11}
\date{2026-09-13}

\begin{document}
\maketitle

\section*{Verdict: FALSE}

We disprove conjecture \texttt{00000000418} as stated. The conjecture is refuted
already at weight $\lambda=(2,1)$, $n=3$, and it is refuted under \emph{each} of
the two readings of ``lattice word'' that the conjecture itself supplies. The
two readings even disagree with each other: the parenthetical
``Yamanouchi (lattice)'' makes the two notions synonyms, whereas the second
clause defines a lattice word as the reading word of a standard tableau, and
reading words are permutations. Since $K_{\lambda,\lambda}=1$ for every
$\lambda$, the conjectured value collapses to $f^\lambda/n!$, while the
proportion of lattice words among Yamanouchi words is identically $1$ under the
synonym reading; and under the strict reading a word has to be a permutation,
so the only weights with nonzero proportions are the one-column rectangles
$(1^n)$.

\section{The conjecture}

Quoted verbatim from \texttt{conjectures/00000000418.md}:

\begin{quote}
\textbf{English.} Definition: A Yamanouchi (lattice) word of weight $\lambda$
is a word in which, in every prefix, the occurrence counts of each letter form
a partition; a lattice word is one appearing as the reading word of a standard
tableau. Conjecture: The proportion of lattice words of weight $\lambda$ among
Yamanouchi words is exactly $f^\lambda/n! \cdot K^{-1}_{\lambda,\lambda}$, and
this ratio attains its maximum $1/2^{n-1}$ when $\lambda$ is a rectangle.
\end{quote}

\begin{quote}\itshape
中文。定义：权 $\lambda$ 的 Yamanouchi(格子)词指每个前缀中字母 $i$
的出现次数构成分拆的词;格子词指出现在某标准表读词中的词。猜想：权
$\lambda$ 的格子词在 Yamanouchi 词中的比例恰为 $f^\lambda/n!\cdot
K^{-1}_{\lambda,\lambda}$,且该比值在 $\lambda$ 取方框形状时达到最大值
$1/2^{n-1}$。
\end{quote}

\section{Conventions}

\begin{definition}[word, content, Yamanouchi word]
A \emph{word} $w=w_1w_2\cdots w_n$ is a finite sequence of positive integers.
Its \emph{content} is the vector $\cont(w)=(c_1,c_2,\dots)$, where $c_i$ is the
number of occurrences of the letter $i$. A \emph{Yamanouchi word of weight
$\lambda$} is a word $w$ with $\cont(w)=\lambda$ such that for every prefix
$w_1\cdots w_m$ $(0\le m\le n)$ the vector of counts
$\bigl(c_1(m),c_2(m),\dots\bigr)$ is a partition, i.e.
$c_1(m)\ge c_2(m)\ge\cdots$.
\end{definition}

Write $n=|\lambda|=\sum_i\lambda_i$, and write $f^\lambda$ for the number of
standard Young tableaux (SYT) of shape $\lambda$ and $K_{\lambda,\mu}$ for the
Kostka number, i.e.\ the number of semistandard Young tableaux (SSYT) of shape
$\lambda$ and content $\mu$.

The classical fact on which both readings turn is the standard bijection
\[
\bigl\{\SYT\text{ of shape }\lambda\bigr\}
\;\longleftrightarrow\;
\bigl\{\text{Yamanouchi words of weight }\lambda\bigr\},
\]
assigning to an SYT $T$ the word $w(T)$ whose $i$-th letter is the index of the
row containing the entry $i$. For example, for $\lambda=(2,1)$ the two SYT
\[
T_1=\begin{array}{|c|c|}\hline 1&2\\\hline 3&\phantom{x}\\\hline\end{array},
\qquad
T_2=\begin{array}{|c|c|}\hline 1&3\\\hline 2&\phantom{x}\\\hline\end{array}
\]
give $w(T_1)=112$ and $w(T_2)=121$. In particular the number of Yamanouchi
words of weight $\lambda$ is $f^\lambda$.

\section{Kostka numbers on the diagonal are all $1$}

\begin{lemma}\label{lem:Kdiag}
For every partition $\lambda$ one has $K_{\lambda,\lambda}=1$. The unique SSYT
of shape $\lambda$ and content $\lambda$ is the \emph{superstandard} tableau,
in which every cell of row $i$ contains the entry $i$.
\end{lemma}

\begin{proof}
Let $T$ be an SSYT of shape $\lambda$ and content $\lambda$. We show by
induction on the row index $i$ that row $i$ of $T$ consists entirely of $i$'s.
For $i=1$: suppose some cell of a row $r\ge 2$ contained a $1$. Row $r$ is
weakly increasing, so every cell of row $r$ weakly to the left of that cell
also contains a $1$; in particular the cell $(r-1,j)$ directly above some cell
$(r,j)$ containing $1$ exists (since $T$ has shape $\lambda$) and satisfies
$T(r-1,j)<T(r,j)=1$ by column strictness, which is impossible because all
entries are positive. Hence all $1$'s lie in the first row. There are exactly
$\lambda_1$ of them, and the first row has exactly $\lambda_1$ cells, so it is
filled with $1$'s. For the induction step, delete the first row; the remaining
array is an SSYT of shape $(\lambda_2,\lambda_3,\dots)$ and content
$(\lambda_2,\lambda_3,\dots)$ with all entries at least $2$; subtracting $1$
from every entry and applying the same argument shows that its first row is
constant equal to $2$, and so on. The superstandard tableau clearly has the
required shape, content and semistandardness, so it is the unique one.
\end{proof}

\begin{corollary}\label{cor:collapse}
The conjectured proportion equals
\[
\frac{f^\lambda}{n!\cdot K_{\lambda,\lambda}}=\frac{f^\lambda}{n!}
\]
for every $\lambda$.
\end{corollary}

\section{Refutation under reading (A): ``Yamanouchi (lattice)'' is a synonym}

\begin{definition}[Reading (A)]
The parenthetical in ``Yamanouchi (lattice) word'' identifies the two terms.
Thus a \emph{lattice word of weight $\lambda$} means a \emph{Yamanouchi word of
weight $\lambda$}, and the set of lattice words of weight $\lambda$ is exactly
the set of Yamanouchi words of weight $\lambda$.
\end{definition}

\begin{theorem}\label{thm:A}
Under reading (A) the proportion of lattice words among Yamanouchi words of
weight $\lambda$ is $1$ for every $\lambda$, while the conjectured value is
$f^\lambda/n!$. At $\lambda=(2,1)$ one has $1\ne 1/3$; hence the conjecture is
false.
\end{theorem}

\begin{proof}
Under reading (A) the two classes of words coincide, so the proportion is
$1$. By Corollary~\ref{cor:collapse} the conjectured value is $f^\lambda/n!$,
which is $<1$ for every partition of size $n\ge 2$ (see
Section~\ref{sec:general}). Evaluating at $\lambda=(2,1)$: $n=3$ and the words
of content $(2,1)$ over $\{1,2\}$ are $112,121,211$. The prefix counts are
\[
\begin{array}{c|ccc|c}
\text{word} & \text{prefix }1 & \text{prefix }2 & \text{prefix }3 & \text{Yamanouchi?}\\
\hline
112 & (1,0) & (2,0) & (2,1) & \text{yes}\\
121 & (1,0) & (1,1) & (2,1) & \text{yes}\\
211 & (0,1) & (1,1) & (2,1) & \text{no, since } (0,1)\text{ is not a partition}
\end{array}
\]
Thus there are exactly $2$ Yamanouchi words of weight $(2,1)$, namely $112$ and
$121$, and the proportion under reading (A) is
\[
\frac{2}{2}=1 .
\]
The hook lengths of the cells of $(2,1)$ are $3,1,1$, so the hook-length
formula gives
\[
f^{(2,1)}=\frac{3!}{3\cdot1\cdot1}=\frac{6}{3}=2,
\]
in agreement with the two SYT $T_1,T_2$ above; moreover $n!=6$ and
$K_{(2,1),(2,1)}=1$ by Lemma~\ref{lem:Kdiag}. The conjectured value is
therefore
\[
\frac{f^{(2,1)}}{3!\cdot K_{(2,1),(2,1)}}=\frac{2}{6\cdot1}=\frac13 .
\]
Since
\[
1\;\ne\;\frac13,
\]
and cross-multiplying gives
\[
2\cdot 6=12\;\ne\;4=2\cdot 2,
\]
the conjecture fails at $\lambda=(2,1)$ under reading (A).
\end{proof}

The enumeration is reproduced by brute force in \texttt{reproduce.py} and is
also formalised in \texttt{lean4/} (the Yamanouchi words of weight $(2,1)$ are
proved to be exactly $112$ and $121$, and the two readings are proved to give
the proportions $2/2$ and $0/2$).

\subsection{The asserted maximum $1/2^{\,n-1}$}

Reading (A) also kills the second clause. At $\lambda=(2,1)$ one has $n=3$, so
the asserted maximum is $1/2^{\,n-1}=1/4$. But the proportion under reading (A)
is $1$, and
\[
1\;>\;\frac14,
\]
so the asserted maximum is not a maximum at all: it is exceeded. Equivalently,
the conjectured value $f^\lambda/n!$ would have to attain $1/2^{\,n-1}$ at a
rectangle, but the proportion itself is $1$ for every shape, and $1\ne1/4$.

\section{Refutation under reading (B): reading words of standard tableaux}

\begin{definition}[Reading (B)]
Read the second clause literally: a \emph{lattice word} is a word that appears
as the reading word of a standard tableau, i.e.\ the word obtained by reading
the entries of some SYT row by row (say from bottom to top, each row left to
right).
\end{definition}

\begin{theorem}\label{thm:B}
Under reading (B) there is no lattice word of weight $(2,1)$, so the proportion
of lattice words among the $2$ Yamanouchi words of weight $(2,1)$ is
\[
\frac{0}{2}=0\;\ne\;\frac13 .
\]
Hence the conjecture is false.
\end{theorem}

\begin{proof}
In a standard tableau the entries $1,2,\dots,n$ each occur exactly once. Hence
every reading word is a permutation of $\{1,\dots,n\}$ and has content
$(1,1,\dots,1)=(1^n)$, so it can have weight $(2,1)$ only if
$(2,1)=(1,1,1)$, which is false. Concretely, the two SYT of shape $(2,1)$ have
reading words $312$ and $213$ (bottom row first), both of content $(1,1,1)$.
Therefore the number of lattice words of weight $(2,1)$ under reading (B) is
$0$, and the proportion is $0/2=0$, whereas the conjectured value is $1/3$ by
Theorem~\ref{thm:A}. This proves the claim.
\end{proof}

\begin{remark}
The two readings are mutually inconsistent already at the level of the
definition: Yamanouchi words of weight $\lambda$ may repeat letters (whenever
some $\lambda_i\ge2$), while reading words of standard tableaux never do. At
$\lambda=(2,1)$ both readings are fatal to the conjecture, and they fail in
opposite directions ($1$ versus $0$, against the claimed $1/3$).
\end{remark}

\section{The general failure: $K_{\lambda,\lambda}=1$ makes the value $f^\lambda/n!$}
\label{sec:general}

By Lemma~\ref{lem:Kdiag}, $K_{\lambda,\lambda}=1$ for every partition, so the
conjectured value is always $f^\lambda/n!$. Under reading (A) the true
proportion is $1$ for every $\lambda$. Since
\[
\sum_{\lambda\vdash n} f^\lambda=n!
\qquad\text{and}\qquad
f^\lambda\ge 1 \text{ for every } \lambda\vdash n,
\]
for $n\ge2$ there are at least two partitions of $n$, hence $f^\lambda<n!$ and
therefore $f^\lambda/n!<1$. So the conjectured proportion $f^\lambda/n!$
differs from the true proportion $1$ for \emph{every} partition $\lambda$ of
every size $n\ge2$; the only case that is not immediately excluded is $n=1$
($\lambda=(1)$, $f^\lambda/n!=1=1$), which is trivial.

\subsection{The maximisers of $f^\lambda/n!$ are not rectangles}

The value $f^\lambda/n!$ is the probability that a uniformly random permutation
of $\{1,\dots,n\}$ has RSK shape $\lambda$. Its maximisers for small $n$ are:
\[
\begin{array}{c|c|c|c}
n & \text{maximising } \lambda & f^\lambda & f^\lambda/n!\\
\hline
3 & (2,1) & 2 & 2/6=1/3\\
5 & (3,1,1) & 6 & 6/120=1/20\\
6 & (3,2,1) & 16 & 16/720=1/45
\end{array}
\]
None of these is a rectangle. At $n=3$ the two rectangle partitions are
$(3)$ and $(1,1,1)$, both with $f^\lambda=1$ and $f^\lambda/n!=1/6<1/3$; the
maximum is attained at the non-rectangle $(2,1)$. So the clause ``attains its
maximum $1/2^{n-1}$ when $\lambda$ is a rectangle'' fails in two independent
ways: the true proportion under reading (A) is $1>1/2^{n-1}$ for every $n\ge2$,
and the conjectured value $f^\lambda/n!$ is not maximised at a rectangle and
does not equal $1/2^{n-1}$ in general (at $n=3$ its maximum is $1/3\ne1/4$).
The case $n=2$ is the only coincidence ($f^{(2)}=f^{(1,1)}=1$, $n!=2$, and
$1/2!=1/2=1/2^{\,n-1}$), but there the true proportion under reading (A) is
still $1\ne1/2$, so the conjecture fails there as well.

\section{Reproduction and formalisation}

The numerical checks are performed by \texttt{reproduce.py} (Python~3,
standard library only), which enumerates all words of content $(2,1)$,
classifies them by the prefix condition, lists the reading words of the two
SYT of shape $(2,1)$, computes $f^{(2,1)}$ by the hook-length formula, computes
$K_{(2,1),(2,1)}$ by brute-force enumeration of semistandard tableaux, asserts
both proportions differ from $1/3$, extends the computation to all partitions
of size $n\le6$ (verifying $\#\text{Yamanouchi}=f^\lambda=\#\SYT$ and
$K_{\lambda,\lambda}=1$ throughout), and prints \texttt{PASS} with exit
status~$0$.

The refutation is formalised in Lean~4 in the accompanying \texttt{lean4/}
project (core Lean only, \texttt{import Std}, no Mathlib, no \texttt{sorry},
toolchain \texttt{leanprover/lean4:v4.33.1}). The formal statements include

\begin{quote}\ttfamily
theorem yamWords21\_eq : yamWords21 = [[1,1,2],[1,2,1]]\\
theorem card\_yamWords21 : yamWords21.length = 2\\
theorem yam\_eq\_latticeReadings : yamWords21 = latticeReadings21\\
theorem numLattice21\_B\_eq : numLattice21\_B = 0\\
theorem f21\_eq : f21 = 2\\
theorem K21\_eq : K21 = 1\\
theorem proportionA\_ne\_claim : propNumA * claimDen \(\ne\) claimNum * propDenA\\
theorem cross\_values : propNumA * claimDen = 12 \(\wedge\) claimNum * propDenA = 4\\
theorem proportionB\_ne\_claim : propNumB * claimDen \(\ne\) claimNum * propDenB\\
theorem quarter\_lt\_one : (1 : Nat) * 1 < 4 * 1\\
theorem proportionA\_exceeds\_max : propNumA * 4 > 1 * propDenA\\
theorem conjecture\_00000000418\_refuted : \dots
\end{quote}

\noindent
Here \texttt{proportionA\_ne\_claim} is the cross-multiplied form of
$2/2\ne2/6$ (namely $12\ne4$) and \texttt{proportionA\_exceeds\_max} is the
cross-multiplied form of $1>1/4$. The document is built with
\texttt{tectonic main.tex}; the resulting PDF is \texttt{build/main.pdf}.

\end{document}
